\documentclass[a4paper,11pt]{ctexart}

\usepackage[margin=2.35cm]{geometry}
\usepackage{amsmath,amssymb,bm}
\usepackage{booktabs}
\usepackage{graphicx}
\usepackage{enumitem}
\usepackage[hidelinks]{hyperref}

\setlist[enumerate]{itemsep=0.25em,topsep=0.4em}
\setlist[itemize]{itemsep=0.25em,topsep=0.4em}
\setlength{\parskip}{0.35em}
\setlength{\parindent}{2em}

\title{\bfseries 光盘光谱仪讲义\\[0.55em]
\large 保护层误差：光栅前面还有一层塑料}
\author{cd\_spectrum}
\date{2026年10月13日}

\begin{document}
\maketitle
\thispagestyle{empty}
\begin{center}
\parbox{0.92\textwidth}{\footnotesize
本讲处理反射式光盘特有的一项误差：数据层埋在塑料下面，入射和衍射都要两次穿过它。
误差分成确定性的折射和随机性的面形、楔角、散射两类。第一类可以靠光路和标定吸收，
第二类只能靠选盘和缩小照明区域来压。与《准直原理》配套阅读；那里的李特罗光路
（DVD、一级、$\alpha=23.4^\circ$、$f=200\,\mathrm{mm}$）是这里的数值基线。}
\end{center}
\vspace{0.6em}
\tableofcontents
\vspace{0.8em}

\section{光到达光栅之前要穿过什么}
\label{sec:stack}

反射式光盘的数据层不在表面上。光从空气进入塑料，到达镀有金属反射层的数据面，
衍射后再穿过同一层塑料出来。各盘的层厚不同。本项目手上的三张盘（合作者记录于
\texttt{已有器材信息.md}）连同蓝光作为参照，列在表~\ref{tab:layers}~中。

\begin{table}[htbp]
\centering\small
\caption{手上的三张盘与蓝光的结构（从读出面往里）}
\label{tab:layers}
\begin{tabular}{llll}
\toprule
盘片 & 层序（读出面 $\to$ 反射层） & 穿过的塑料 & 轨距 \\
\midrule
CD-R（竹子） & PC 衬底 $\to$ 染料 $\to$ 银/金 & $\approx 1.2\,\mathrm{mm}$ & $1.6\,\mu\mathrm{m}$ \\
DVD-R（love） & 两个 $0.6\,\mathrm{mm}$ 半盘粘合 & $\approx 0.6\,\mathrm{mm}$ & $0.74\,\mu\mathrm{m}$ \\
DVD+R（深蓝） & 同上 & $\approx 0.6\,\mathrm{mm}$ & $0.74\,\mu\mathrm{m}$ \\
蓝光（参照） & 覆盖层 $\to$ 记录层 & $\approx 0.1\,\mathrm{mm}$ & $0.32\,\mu\mathrm{m}$ \\
\bottomrule
\end{tabular}
\end{table}

标签上的「竹子 / love / 深蓝」只是盘面图案，不是制造商。每张盘中心孔附近的注塑码区
压着真实的厂家码（如 Ritek、CMC、MCC），测量档案里应先抄下这组码。还要确认盘是
空白还是已写：空白可写盘的预刻槽连续光滑，写过的盘槽被记录标记打断，衍射效率
下降。

聚碳酸酯在可见光波段的折射率约为
\[
n \approx 1.58,\qquad
n(486\,\mathrm{nm})-n(656\,\mathrm{nm})\approx 0.01,
\]
即全可见光段相对变化约 $0.6\%$。

\textbf{染料层在衬底与金属层之间}，反射式用法里光要两次穿过它。染料按波段吸收：
菁染料（盘面偏蓝绿）在红段吸收更强，酞菁染料（浅金色）更透明，偶氮染料呈深蓝色
——「深蓝」很可能就是偶氮系。所以三张盘比较亮度时，差别里有一部分是染料吸收，
不是光栅效率，两者必须分开报。

\textbf{DVD-R 与 DVD+R 在光学上几乎相同}：轨距、层厚一致，差别只在寻址方式
（DVD-R 的岸台预刻坑对 DVD+R 的槽摆动 ADIP），槽的径向调制只有纳米量级。
因此「love」和「深蓝」是天然的同结构重复性对照对：它们之间测出的差别是品牌与
工艺差异，不是标准差异。CD-R（竹子）则是离群点：层厚两倍、轨距两倍，预期谱线
最宽。

这层塑料引起的后果要分成两类。

\textbf{确定性效应}：平板折射改变光在介质内的角度，楔角造成整体偏折。这两项
对固定的盘、固定的照明区域是可重复的，标定可以吸收。

\textbf{随机性缺陷}：面形不平、局部翘曲、体散射、应力双折射、前表面反射的鬼影。
这些随盘片和照明位置变化，不能标定，只能在结构上压制。

\section{平行板对平行光不产生像差}
\label{sec:planewave}

先确立一个事实：严格的平面波（平行光）通过严格的平面平行平板，出射后仍然是
严格的平面波，方向按折射定律精确地折转，只是整个光束横向平移
\begin{equation}
\label{eq:shift}
\Delta = t\sin\alpha\left(1-\frac{\cos\alpha}{\sqrt{n^2-\sin^2\alpha}}\right).
\end{equation}
不产生任何像差。

这正是《准直原理》里光路的好处：光盘前面是准直光，保护层在「平行光过平行板」
的工况下，像差的主项在结构上已经被消掉。这也是为什么准直必须放在盘之前，
而不是用会聚光直接打在盘面上。

\section{光束不平行时：焦点偏移与球差}
\label{sec:focus}

如果用会聚或发散的光穿过这层塑料（相机镜头凑近盘面直接拍、眼睛贴近看、
把成像透镜放在保护层之前），平面板在会聚光束中会产生像差。

首先是整个焦点沿光轴后移
\begin{equation}
\label{eq:fshift}
\Delta z = t\left(1-\frac{1}{n}\right).
\end{equation}
取 $n=1.58$：
\[
\begin{array}{ll}
\text{CD：} & 1.2\times(1-1/1.58) \approx 0.44\,\mathrm{mm},\\[0.2em]
\text{DVD：} & 0.6\times(1-1/1.58) \approx 0.22\,\mathrm{mm},\\[0.2em]
\text{蓝光：} & 0.1\times(1-1/1.58) \approx 0.04\,\mathrm{mm}.
\end{array}
\]
对照《准直原理》的离焦容许（目标 $0.6\,\mathrm{mm}$，上限 $1\,\mathrm{mm}$）：
CD 的焦点偏移已经占到容许的一半以上。在焦点移动之外，会聚光过平板还产生球差，
大小随层厚 $t$ 和光束口径的三次方增长。纸盒目视仪能用，是因为它对分辨率的要求
低；高分辨率光路里，会聚必须放在保护层之外。

\section{折射对光栅方程的影响}
\label{sec:refraction}

衍射发生在塑料内部。光栅方程里的角度应该是介质内的角度
$\alpha'$、$\beta'$：
\begin{equation}
\label{eq:grating}
m\lambda = d\,(\sin\alpha' + \sin\beta'),\qquad
\sin\alpha' = \frac{\sin\alpha}{n},\quad
\sin\beta' = \frac{\sin\beta}{n}.
\end{equation}
如果直接拿空气里的角度代入真空光栅方程，得到的波长刻度是错的。

以李特罗条件为例。空气里 DVD 对 $589\,\mathrm{nm}$ 的李特罗角是
$\alpha=23.4^\circ$；塑料内部实际发生衍射的角度是
\begin{equation}
\sin\alpha' = \frac{\sin 23.4^\circ}{1.58} = 0.252,
\qquad
\alpha' = 14.6^\circ.
\end{equation}
但由于入射角、衍射角同时被同一个 $n$ 缩放，整条谱线的刻度变化是\textbf{系统性的}：
对固定的盘、固定的几何，像素到波长的映射只是换了一个确定的函数。用已知谱线
做多点标定之后，这一项大部分被拟合吸收。

吸收不了的是残余的色散项：$n$ 随波长变化约 $0.6\%$，所以介质内的角度关系
随波长有轻微不同的缩放。若只用两点直线标定或简单的正弦拟合，可见光全段会留下
纳米量级的非线性残差。用三条以上谱线做多项式标定，这一项还能再压一个量级。

\textbf{结论}：折射本身不是障碍，不标定才是。

\section{楔角与面形}
\label{sec:wedge}

盘面不严格平行时，光束被偏折
\begin{equation}
\delta_{\mathrm{dev}} = (n-1)\,\delta_{\mathrm{wedge}} \approx 0.58\,\delta_{\mathrm{wedge}},
\end{equation}
而且偏折量随 $n(\lambda)$ 随波长变化，造成谱线位置随波长的系统性错位。
这一项确定、可重复，同一装夹下标定可以吸收；代价是重新装盘、换一个照明位置，
楔角变了，标定就漂移。

面形误差（局部不平、翘曲）更严重：平行光进去，出来不再平行。按《准直原理》的
换算 $\Delta\lambda=(d\cos\alpha/m)\,\Delta\alpha$，面形引起的角度散布直接变成
谱线展宽，而且不可标定。能做的只有两件事：
\begin{itemize}
\item 掩模只留几毫米宽的窄带，把照明限制在局部平整的区域内；
\item 同一张盘上多试几个位置，选谱线最窄的那一处，并记录取样位置。
\end{itemize}

\section{体散射、双折射与鬼影}
\label{sec:scatter}

\textbf{体散射与双折射。} 注塑聚碳酸酯内部有雾度和应力双折射，光两次穿过后，
一部分变成杂散背景。后果是谱线对比度下降：钠双线中间的谷被填平。这一项既不能
靠光路消，也不能靠标定消，只能选盘。

\textbf{前表面鬼影。} 空气--塑料界面按菲涅尔公式反射约 $4\%$--$5\%$ 的光。
这束光不经过数据层，不产生色散，但双程平移之后进入探测器，在谱旁边形成
位置固定的鬼线，容易被误认成真实谱线。处理办法：
\begin{itemize}
\item 安排几何，让前表面反射光落到掩模和探测器视场之外；
\item 挡光：狭缝、光阑和遮光罩把不参与衍射的反射挡掉；
\item 判据：鬼线的位置不随波长标定移动成谱线该有的样子，且遮住数据层区域
（只留前表面）时它仍在。
\end{itemize}

\section{光路设计能消掉多少}
\label{sec:mitigation}

\begin{table}[htbp]
\centering
\caption{保护层各项误差的对策}
\label{tab:mitigation}
\begin{tabular}{lll}
\toprule
误差项 & 性质 & 对策 \\
\midrule
平行板折射 & 确定性 & 平行光通过即无像差；标定吸收刻度变化 \\
塑料色散 $n(\lambda)$ & 确定性 & 三条以上谱线做多项式标定 \\
楔角偏折 & 半确定性 & 固定装夹；换位置后重新标定 \\
会聚光中的焦移与球差 & 结构性 & 准直放在盘之前，会聚光不穿保护层 \\
面形与翘曲 & 随机 & 窄带掩模；盘上选点 \\
体散射与双折射 & 随机 & 只能选盘；或用蓝光（层薄 $12$ 倍） \\
前表面鬼影 & 结构性 & 几何避让 + 挡光 \\
\bottomrule
\end{tabular}
\end{table}

选蓝光时覆盖层只有约 $0.1\,\mathrm{mm}$，面形、焦移、散射各项天然比 CD 小
一个量级，代价是轨距 $320\,\mathrm{nm}$ 必须用大入射角才能让可见光一级衍射有解，
整套几何要重新设计。撕掉衬底改透射式可以彻底绕开保护层，但那是另一套仪器，
不属于本讲的反射式路线。

\section{不考虑保护层，输出是什么样}
\label{sec:ignored}

\begin{enumerate}
\item \textbf{波长刻度。} 用空气角代入真空光栅方程，刻度整体错误；标定后大部分
被拟合吸收，剩下塑料色散造成的非线性残差，可见光全段纳米量级。
\item \textbf{谱线形状。} 半高宽比缝宽和焦距预言的更大，且找不到是哪个环节宽的。
会聚光穿保护层时，CD 的焦移约 $0.44\,\mathrm{mm}$，叠加球差。
\item \textbf{假谱线。} 前表面鬼影形成位置固定的卫星线，容易被误认。
\item \textbf{对比度。} 体散射抬高背景，钠双线的谷变浅甚至填平，实测分辨本领
低于理论值，且无法归因。
\item \textbf{重复性。} 重新装盘后楔角和面形变了，标定漂移，同一条线的多次
测量不可重复。
\end{enumerate}

这些症状里，「找得到规律」的部分（1、3 的位置关系）可以靠标定和几何处理；
「找不到规律」的部分（2、4、5）就是保护层留给反射式光盘的下限，要靠第
\ref{sec:experiment}~节的实验定量。

\section{测量方案}
\label{sec:experiment}

把保护层当成一个待测变量，而不是先验假设。原则：每次只换一个量。

三张盘进光路之前，先各建一行档案：内圈厂家码、总厚（千分尺，三个半径乘四个
方位角）、衬底光学厚度、折射率、楔角、偏振片照片、记录面颜色、空白与否。
具体的测量方法与精度分析见仓库文档《三张盘的结构分析与测量方案》。

\begin{enumerate}
\item \textbf{换盘实验。} 缝宽、焦距、入射角、光源、相机设置全部不变，只把
CD-R（竹子）、DVD-R（love）、DVD+R（深蓝）轮换装到同一位置。对同一条窄线
（低压钠灯 $589\,\mathrm{nm}$ 或汞 $546\,\mathrm{nm}$）测半高宽。预期：
两张 DVD 接近，CD-R 最宽；若两张 DVD 差别明显，归因于品牌工艺而非标准。
\item \textbf{染料吸收。} 比较三张盘同一谱线的相对强度随波长的变化，
把染料滤光从光栅效率里分离出来。
\item \textbf{换位置实验。} 同一张盘，移动掩模照到不同半径、不同方位角，记录
谱线半高宽和弯曲的变化，给出面形误差的分布。
\item \textbf{鬼影检查。} 遮住数据层只留前表面反射，看鬼线是否还在、在什么位置，
确认它已被挡在视场之外。
\item \textbf{装夹重复性。} 同一位置重复装夹数次，比较标定系数和谱线半高宽的
漂移，给出楔角项的实际大小。
\end{enumerate}

所有实验保存原始图与暗场图，记录盘片型号、取样位置、缝宽、角度与相机参数。
报告分辨本领时，把这些数据作为谱线宽度能否归因的依据。

\section{小结}

反射式光盘的光栅前面有一层 $0.1$--$1.2\,\mathrm{mm}$ 的塑料。它对严格平行光
不产生像差，所以准直放在盘之前就把主项消掉了。折射和塑料色散是确定性的，
多点标定可以吸收到亚纳米量级。剩下的是随机的面形、散射与鬼影：面形靠窄带
掩模和选点压制，散射只能靠选盘（蓝光层最薄），鬼影靠几何避让和挡光。

不处理它的代价是：谱线比预言的宽且无法归因、出现假谱线、双线谷被填平、
装夹之间不可重复。保护层因此是反射式光盘分辨本领的下限来源之一，
必须用换盘实验定量，不能只靠计算。手上的三张盘里，CD-R（竹子）层最厚、
风险最大；DVD-R（love）与 DVD+R（深蓝）结构相同，互为重复性对照。
每张盘的厚度、折射率、楔角先按配套文档测好存档，换盘实验的谱线差别
才有归因的依据。

\end{document}
